
%(BEGIN_QUESTION)
% Copyright 2013, Tony R. Kuphaldt, released under the Creative Commons Attribution License (v 1.0)
% This means you may do almost anything with this work of mine, so long as you give me proper credit

\noindent
Sammenlign og diskuter forskjellene mellom {\it laminær} og {\it turbulent} strømning med hensyn til følgende kriterier:

\vskip 10pt

\begin{itemize}
\item{} Hvilket regime skaper minst {\it strømningsmotstand} (friksjonstap) gjennom en lang rørlengde? 
\vskip 10pt
\item{} Hvilket regime er best for å blande væsker sammen i et rørsystem?
\vskip 10pt
\item{} Hvilket regime er best for å sikre en grundig reaksjon mellom kjemiske reaktanter i et rørsystem?
\vskip 10pt
\item{} Hvilket regime foretrekkes inne i en varmeveksler, for å sikre maksimal varmeoverføring?
\end{itemize}

\vskip 30pt

\noindent
Undersøk følgende formler for å beregne Reynolds-tall:

\vskip 10pt

For å beregne Reynolds-tall med engelske enheter (gassstrømning):

$$\hbox{Re} = {{(6.32) \rho Q} \over {D \mu}}$$

\noindent
Hvor,

Re = Reynolds-tall (enhetsløst)

$\rho$ = Massetetthet for gass, i pund (masse) per kubikkfot (lbm/ft$^{3}$)

$Q$ = Strømningshastighet, standard kubikkfot per time (SCFH)

$D$ = Rørets diameter, i tommer (in)

$\mu$ = Absolutt viskositet for gassen, i centipoise (cP)

\vskip 20pt \goodbreak

For å beregne Reynolds-tall med relativ tetthet i stedet for massetetthet (væskestrømning):

$$\hbox{Re} = {{(3160) G_f Q} \over {D \mu}}$$

\noindent
Hvor,

Re = Reynolds-tall (enhetsløst)

$G_f$ = Relativ tetthet for væsken (enhetsløs)

$Q$ = Strømningshastighet, gallons per minutt (GPM)

$D$ = Rørets diameter, i tommer (in)

$\mu$ = Absolutt viskositet for væsken, i centipoise (cP)

\vskip 20pt

Identifiser nå kvalitativt hvilken retning hver variabel i formelen må endres (f.eks.\ {\it øke} i motsetning til å {\it avta}) for å fremme turbulens i en strøm, forutsatt at alle andre faktorer forblir uendret.

\vskip 20pt \vbox{\hrule \hbox{\strut \vrule{} {\bf Forslag til sokratisk diskusjon} \vrule} \hrule}

\begin{itemize}
\item{} For en gitt {\it strømningshastighet} i væsken vil Reynolds-tallet minke hvis rørdiameteren minker. Med dette i bakhodet: forklar hvorfor vi ser $D$ i nevneren i disse brøkene, og ikke i telleren.
\end{itemize}

\underbar{file i01299}
%(END_QUESTION)





%(BEGIN_ANSWER)

\begin{itemize}
\item{} Hvilket regime skaper minst {\it strømningsmotstand} (friksjonstap) gjennom en lang rørlengde?  {\it Laminær}
\vskip 5pt
\item{} Hvilket regime er best for å blande væsker sammen i et rørsystem?  {\it Turbulent}
\vskip 5pt
\item{} Hvilket regime er best for å sikre en grundig reaksjon mellom kjemiske reaktanter i et rørsystem?  {\it Turbulent}
\vskip 5pt
\item{} Hvilket regime foretrekkes inne i en varmeveksler, for å sikre maksimal varmeoverføring?  {\it Turbulent}
\end{itemize}

\vskip 10pt

\noindent
Når alle andre faktorer er like, fremmes turbulens av følgende:

\begin{itemize}
\item{} Økt strømningshastighet ($Q$)
\vskip 5pt
\item{} Økt væsketetthet ($\rho$ eller $G_f$)
\vskip 5pt
\item{} Mindre rørdiameter ($D$)
\vskip 5pt
\item{} Lavere viskositet ($\mu$)
\end{itemize}



%(END_ANSWER)





%(BEGIN_NOTES)

\noindent
Andre formler for Reynolds-tall:

\vskip 10pt

For å beregne Reynolds-tall med metriske enheter:

$$\hbox{Re} = {{D \overline{V} \rho} \over \mu}$$

\noindent
Hvor,

Re = Reynolds-tall (enhetsløst)

$D$ = Rørets diameter, i meter (m)

$\overline{V}$ = Gjennomsnittlig hastighet for væske, i meter per sekund (m/s)

$\rho$ = Massetetthet for væske, i kilo per kubikkmeter (kg/m$^{3}$)

$\mu$ = Absolutt viskositet for væske, i Pascal-sekunder (Pa $\cdot$ s)

\vskip 60pt \goodbreak

For å beregne Reynolds-tall med engelske enheter (væskestrømning):

$$\hbox{Re} = {{(50.7) \rho Q} \over {D \mu}}$$

\noindent
Hvor,

Re = Reynolds-tall (enhetsløst)

$\rho$ = Massetetthet for væske, i pund (masse) per kubikkfot (lbm/ft$^{3}$)

$Q$ = Strømningshastighet, gallons per minutt (GPM)

$D$ = Rørets diameter, i tommer (in)

$\mu$ = Absolutt viskositet for væsken, i centipoise (cP)


%INDEX% Physics, dynamic fluids: laminar versus turbulent flow regimes
%INDEX% Physics, dynamic fluids: Reynolds number

%(END_NOTES)
